% =================================================================
% CHƯƠNG 3: THỰC NGHIỆM VÀ KẾT QUẢ
% =================================================================
\chapter{THỰC NGHIỆM VÀ KẾT QUẢ}
\label{ch:thuc_nghiem_ket_qua}

Chương này trình bày chi tiết quy trình thiết lập thực nghiệm, mô tả bộ dữ liệu thương mại điện tử Olist, các bước tiền xử lý lọc chuỗi thưa và gom cụm địa lý. Tiếp theo, nghiên cứu định nghĩa các độ đo đánh giá hiệu năng dự báo nhu cầu vận chuyển. Cuối cùng, kết quả đối sánh dự báo chuỗi thời gian giữa mô hình đề xuất Mamba4Cast với các baseline và kết quả tối ưu hóa mạng lưới kho trung chuyển đa thời kỳ (CWLP) bằng hai giải thuật MILP và Benders Decomposition sẽ được trình bày và thảo luận chi tiết.

% =================================================================
% 3.1 QUY TRÌNH THỰC NGHIỆM
% =================================================================
\section{Quy trình thực nghiệm}
\label{sec:quy_trinh_thuc_nghiem}

Đồ án ứng dụng mô hình tích hợp theo dạng \textit{Forecast-Then-Optimize} gồm 2 giai đoạn: Giai đoạn 1 tiến hành dự b
<truncated 14811 bytes>
 điểm hội tụ hoàn toàn.
    \item \textbf{Đánh giá thời gian chạy:} MILP giải trực tiếp bằng CBC chỉ mất trung bình $0{,}92$ giây, trong khi Benders Decomposition mất tới $31{,}03$ giây. Chi phí thời gian của Benders lớn hơn do overhead giao tiếp và giải lặp đi lặp lại bài toán Master và bài toán đối ngẫu con (Dual Subproblem) thông qua giao diện lập trình PuLP trong Python.
    \item \textbf{Nghiệm vị trí kho mở tối ưu:} Cả hai giải thuật đều cho thấy sự đồng thuận cao trong việc đề xuất thiết lập mạng lưới kho bán tĩnh (semi-static), mở $4$ kho trung chuyển tại các bang kinh tế trọng điểm của Brazil: nội thành và ngoại ô São Paulo (\texttt{04}/\texttt{05}, \texttt{08}, \texttt{13}) và bang Minas Gerais (\texttt{35}), khớp hoàn toàn với bản đồ phân bổ nhu cầu tiêu dùng tập trung của Brazil.
\end{itemize}
